\documentclass[11pt]{article}
\usepackage[margin=1in]{geometry}
\usepackage{amsmath,amssymb,amsthm}
\usepackage[T1]{fontenc}
\usepackage{lmodern}
\usepackage{xcolor}
\usepackage{booktabs}
\usepackage{longtable}
\usepackage{array}
\usepackage[hidelinks]{hyperref}
\usepackage{url}

\newtheorem{theorem}{Theorem}
\newtheorem{proposition}[theorem]{Proposition}
\newtheorem{lemma}[theorem]{Lemma}
\newtheorem{corollary}[theorem]{Corollary}
\theoremstyle{definition}
\newtheorem{definition}[theorem]{Definition}
\theoremstyle{remark}
\newtheorem{remark}[theorem]{Remark}

% Code-like text (not checked by a name gate).
\newcommand{\lean}[1]{\texttt{#1}}
% A declaration of one of the two developments; every \decl{...} is checked by namegate_c.sh
% (it must be in that repository's gate.sh REQUIRED list).
\DeclareUrlCommand\decl{\urlstyle{tt}\def\UrlBreaks{\do\_\do\.}}
% A Mathlib declaration (pinned revision 0df444a3); not one of ours.
\DeclareUrlCommand\mdecl{\urlstyle{tt}\def\UrlBreaks{\do\_\do\.}}
\newcommand{\RR}{\mathbb{R}}
\newcommand{\CC}{\mathbb{C}}
\newcommand{\PH}[1]{\textcolor{red}{\textbf{[PLACEHOLDER:} #1\textbf{]}}}

% ---------------------------------------------------------------------------------------------
% PINS. Both developments are published; the \src anchors below link to these commits.
% ---------------------------------------------------------------------------------------------
\newcommand{\commitof}[1]{\csname commitof:#1\endcsname}
\expandafter\def\csname commitof:barbier\endcsname{7baae0430954ce78ddf8d6e6bf5a4b841d7b2367}
\expandafter\def\csname commitof:fourvertex\endcsname{14d10449dd1a206f5342099d65d5f4180830d423}
\newcommand{\repo}[1]{https://github.com/jpbald93/#1}
\newcommand{\src}[4]{\ifcsname commitof:#1\endcsname
  \href{\repo{#1}/blob/\commitof{#1}/#2/#3.lean\#L#4}{\texttt{#3}:#4}\else\texttt{#3}:#4\fi}

\title{Barbier's theorem for the Reuleaux triangle and the convex four-vertex theorem\\
in Lean~4, through support functions}
\author{Josh Bald\\ \small Independent researcher, Ontario, Canada\\ \small\texttt{jpbald93@gmail.com}\\
\small ORCID: \texttt{0009-0002-1317-6489}
\thanks{\copyright~2026 Josh Bald. Licensed under CC BY 4.0 (\url{https://creativecommons.org/licenses/by/4.0/}). The accompanying Lean code is licensed under the Apache License 2.0.}}
\date{10 October 2026}

\emergencystretch=3em
\begin{document}
\sloppy
\maketitle

\begin{abstract}
We give machine-checked proofs, in Lean~4 with Mathlib, of two classical theorems on convex plane
curves, both developed in one model: a curve is described by its support function $h$, and is
traced as $\gamma_h(\theta) = (h(\theta) + i\,h'(\theta))\,e^{i\theta}$.
The first is Barbier's theorem: a curve of constant width $w$ has perimeter $\pi w$. We prove it
not only for twice continuously differentiable support functions but for support functions that are
merely $C^1$, with $h'$ differentiable and $h+h''\ge 0$ outside a countable set and $h+h''$
integrable there. This covers the Reuleaux
triangle, which we construct, show to satisfy the hypotheses, and show not to be twice
continuously differentiable, so the smooth theorem does not reach it. For the Reuleaux triangle we also identify the perimeter
with Mathlib's length and exhibit it as a Mathlib convex body of constant width. The second is the four-vertex theorem for
ovals given by $C^3$ support functions with $h+h''>0$: the curvature is constant or has at least four
alternating local extrema in a period. We tie each of our own definitions to a standard notion: the perimeter to
Mathlib's length (total variation) --- for smooth support functions, and separately for the Reuleaux
triangle --- via a lemma, not in Mathlib, that the length of a $C^1$ curve
is finite and equals $\int\lvert\gamma'\rvert$; the support function to an actual convex set, whose
support values and widths we compute; and our curvature to the textbook formula
$(x'y''-y'x'')/\lvert\gamma'\rvert^3$, for a curve we prove to be simple and closed.
All proofs use only Lean's three standard axioms. We know of no previous complete formal proof of
either theorem in any proof assistant, and no formalisation of the Reuleaux triangle.
\end{abstract}

\section{Introduction}

A compact convex set in the plane has \emph{constant width} $w$ if every pair of parallel support
lines is at distance $w$. Barbier's theorem \cite{barbier1860} says that every such set has perimeter $\pi w$:
the circle of diameter $w$ is far from the only example, but all of them have the circle's
perimeter. The best-known non-circular example is the Reuleaux triangle, bounded by three circular
arcs of radius $w$, each centred at the opposite corner; see \cite{martini2019} for the theory of
constant width.

The four-vertex theorem \cite{mukhopadhyaya1909,kneser1912} says that the
curvature of a simple closed smooth plane curve is either constant or has at least four local
extrema. A standard proof for strictly convex curves goes through the support function and a
Sturm--Hurwitz-type lemma: a continuous $2\pi$-periodic function orthogonal to $1$, $\cos$ and $\sin$
changes sign at least four times in a period, unless it vanishes; see \cite{osserman1985,gluck1971,deturck2007} for the history and the converse.

Both theorems appear on the ``1000+ theorems'' list that Mathlib tracks, and neither has a formal
proof there (Section~\ref{sec:related}). This paper gives complete Lean~4 proofs of both, in a
common model, and says exactly how far each statement reaches.

\paragraph{Contributions.}
\begin{enumerate}
\item \textbf{Barbier, with corners} (Theorem~\ref{thm:barbierpw}): a curve whose support function is
$C^1$, $2\pi$-periodic and of constant width $w$, with $h'$ differentiable and $h+h''\ge 0$ outside a
countable set and $h+h''$ integrable (Definition~\ref{def:pcw}), has
perimeter $\pi w$. The smooth version (Theorem~\ref{thm:barbier}) is a special case
(Proposition~\ref{prop:topw}).
\item \textbf{The Reuleaux triangle} (Theorem~\ref{thm:reuleaux}): an explicit support function,
proved to satisfy the hypotheses of Theorem~\ref{thm:barbierpw}, to have perimeter $\pi w$, to have
three corners pairwise at distance $w$ and sides that are arcs of radius $w$ about the opposite
corner, and \emph{not} to be twice continuously differentiable; and, with Mathlib's length in
place of our own perimeter, that the length is $\pi w$ (Theorem~\ref{thm:reuleaux}(f)).
\item \textbf{The four-vertex theorem for ovals given by $C^3$ support functions with $h+h''>0$}
(Theorems~\ref{thm:fourvertex} and~\ref{thm:fvtextbook}), with the Sturm--Hurwitz lemma
(Lemma~\ref{lem:sturm}).
\item \textbf{Bridges to standard notions} (Section~\ref{sec:bridges}): the length of a $C^1$ curve
is finite and equals $\int\lvert\gamma'\rvert$ (Theorem~\ref{thm:length}); our perimeter is
Mathlib's length for smooth support functions (Corollary~\ref{cor:perlen}) and, separately, for the
Reuleaux triangle (Theorem~\ref{thm:reuleaux}(f)); the support function $h$ is the support function of
an explicit convex set of constant width $w$, packaged as a Mathlib convex body
(Theorems~\ref{thm:body} and~\ref{thm:convbody}); our curvature is the
textbook signed curvature of a simple closed curve (Theorem~\ref{thm:curv}).
\end{enumerate}

\paragraph{What is not claimed.} We do not prove that every convex body of constant width has a
support function of the form we assume, so Theorem~\ref{thm:barbierpw} is Barbier's theorem for
curves \emph{given} by such support functions, not for arbitrary convex bodies. The four-vertex
theorem is proved for ovals given by $C^3$ support functions with $h+h''>0$, not for general
simple closed curves. Section~\ref{sec:limits} lists these and the other limitations.

\subsection{Availability}\label{sec:avail}
The proofs are in two Lean~4 developments, \lean{barbier} (ten files, 2{,}746 lines) and
\lean{fourvertex} (six files, 858 lines), on Lean v4.33.1 \cite{lean4} with Mathlib \cite{mathlib2020} pinned at revision
\texttt{0df444a3}. They are published as \url{https://github.com/jpbald93/barbier} (commit
\texttt{7baae04}) and \url{https://github.com/jpbald93/fourvertex} (commit \texttt{14d1044}); every
source reference in this paper links to those commits. Each development has a gate script: it rejects \lean{sorry}, \lean{admit},
\lean{native\_decide}, \lean{axiom} and meta-programming commands in the sources, builds the library,
and checks with \lean{\#print axioms} that every listed declaration (62 in \lean{barbier}, 48 in
\lean{fourvertex}) depends only on \lean{propext}, \lean{Classical.choice} and \lean{Quot.sound}.
Every declaration named in this paper is in those lists (Section~\ref{sec:verify}).

\section{The support-function model}\label{sec:model}

We identify the plane with $\CC$. For $h:\RR\to\RR$ write $h'$ for its derivative, with the usual
Lean convention that the derivative is $0$ where it does not exist.

\begin{definition}[curve and perimeter]\label{def:curve}
The \emph{curve} of $h$ is $\gamma_h(\theta) = (h(\theta) + i\,h'(\theta))\,e^{i\theta}$, and its
\emph{perimeter} is $P(h) = \int_0^{2\pi}\lvert\gamma_h'(\theta)\rvert\,d\theta$
(\decl{Barbier.curve}, \src{barbier}{Barbier}{Basic}{28}; \decl{Barbier.perimeter}, \src{barbier}{Barbier}{Basic}{31}).
\end{definition}

The geometry behind the formula: $\langle\gamma_h(\theta), e^{i\theta}\rangle = h(\theta)$, so
$\gamma_h(\theta)$ lies on the line with outer normal $e^{i\theta}$ at signed distance $h(\theta)$
from the origin, and when $h$ is $C^2$ the velocity is
$\gamma_h'(\theta) = (h+h'')(\theta)\cdot i e^{i\theta}$, tangent to that line
(\decl{Barbier.bridge_curve}, \src{barbier}{Barbier}{Bridge}{37}). So $\rho = h+h''$ is the radius
of curvature, and the curve is convex where $\rho\ge 0$.

\begin{definition}[smooth constant-width support function]\label{def:cws}
$h$ is a \emph{smooth constant-width support function of width $w$}
(\decl{Barbier.ConstWidthSupport}, \src{barbier}{Barbier}{Basic}{21}) if $h$ is $C^2$ and
$2\pi$-periodic, $h(\theta)+h(\theta+\pi) = w$ for all $\theta$, and $h+h''\ge 0$ everywhere.
\end{definition}

\begin{definition}[piecewise constant-width support function]\label{def:pcw}
$h$ is a \emph{piecewise constant-width support function of width $w$}
(\decl{Barbier.PiecewiseConstWidth}, \src{barbier}{Barbier}{Piecewise}{23}) if $h$ is $C^1$ and
$2\pi$-periodic, $h(\theta)+h(\theta+\pi) = w$ for all $\theta$, there is a countable set $S$ such
that for every $\theta\notin S$ the function $h'$ is differentiable at $\theta$ and
$h(\theta)+h''(\theta)\ge 0$, and $h+h''$ is integrable on $[0,2\pi]$.
\end{definition}

Countability of $S$ is exactly what the fundamental theorem of calculus with a countable exceptional
set needs (Mathlib's \mdecl{intervalIntegral.integral_eq_of_hasDerivAt_off_countable_of_le}). It
allows finitely many corners per period, as in the Reuleaux triangle, and more.

\section{Barbier's theorem}\label{sec:barbier}

\begin{theorem}[Barbier, smooth]\label{thm:barbier}
If $h$ is a smooth constant-width support function of width $w$, then $P(h) = \pi w$
(\decl{Barbier.barbier}, \src{barbier}{Barbier}{Basic}{110}).
\end{theorem}

\begin{proof}[Proof sketch]
Since $h+h''\ge 0$, $\lvert\gamma_h'\rvert = h+h''$ (\decl{Barbier.norm_deriv_curve_of_convex}, \src{barbier}{Barbier}{Basic}{80}). Periodicity of $h'$ gives $\int_0^{2\pi}h'' = 0$, so
$P(h) = \int_0^{2\pi}h$ (\decl{Barbier.perimeter_eq_integral}, \src{barbier}{Barbier}{Basic}{99}).
Splitting $[0,2\pi]$ at $\pi$ and using $h(\theta)+h(\theta+\pi) = w$ gives $\int_0^{2\pi}h = \pi w$
(\decl{Barbier.integral_support_eq}, \src{barbier}{Barbier}{Basic}{34}); this step needs only
continuity of $h$.
\end{proof}

\begin{proposition}\label{prop:topw}
Every smooth constant-width support function of width $w$ is a piecewise one of width $w$
(\decl{Barbier.ConstWidthSupport.toPiecewise}, \src{barbier}{Barbier}{Piecewise}{57}).
\end{proposition}

\begin{theorem}[Barbier, with corners]\label{thm:barbierpw}
If $h$ is a piecewise constant-width support function of width $w$, then $P(h) = \pi w$
(\decl{Barbier.barbier_piecewise}, \src{barbier}{Barbier}{Piecewise}{72}).
\end{theorem}

\begin{proof}[Proof sketch]
Off $S$ the derivative of $\gamma_h$ is still $(h+h'')\,i e^{i\theta}$, which needs only $h$ to be
$C^1$ and $h'$ differentiable at that point, so $\lvert\gamma_h'\rvert = h+h''$ off $S$. A countable
set is Lebesgue-null, so $P(h) = \int_0^{2\pi}(h+h'')$. The function $h'$ is continuous, periodic,
and differentiable off $S$ with integrable derivative, so the countable-exception fundamental theorem
of calculus gives $\int_0^{2\pi}h'' = h'(2\pi)-h'(0) = 0$. The last step is as before.
\end{proof}

The care needed is in the junk values: where a derivative does not exist, Lean assigns it the value
$0$, so on $S$ both $h''$ and $\gamma_h'$ may take junk values, and the proof never uses those values
except through null sets.

\paragraph{Examples.} The circle of diameter $w$ (\decl{Barbier.circle_perimeter}, \src{barbier}{Barbier}{Examples}{23}) and the smooth family
$h(\theta) = a + b\cos 3\theta$ with $0\le 8b\le a$, non-circular when $b>0$ (\decl{Barbier.trefoil_perimeter}, \src{barbier}{Barbier}{Examples}{66}) are instances of Theorem~\ref{thm:barbier}. For the circle, the
curve is Mathlib's \mdecl{circleMap} and the perimeter is $2\pi\lvert r\rvert$
(\decl{Barbier.bridge_perimeter}, \src{barbier}{Barbier}{Bridge}{122}).

\section{The Reuleaux triangle}\label{sec:reuleaux}

\begin{definition}[Reuleaux support function]\label{def:reuleaux}
For $w\in\RR$ let
\[
  g_w(p) = \begin{cases} \dfrac{w}{\sqrt3}\cos p, & p\le \pi/6,\\[1ex]
  w - \dfrac{w}{\sqrt3}\cos(p-\pi/3), & p > \pi/6,\end{cases}
\]
and let $R_w(\theta) = g_w(p)$, where $p$ is the unique number in $[-\pi/6, -\pi/6+2\pi/3)$
congruent to $\theta$ modulo $2\pi/3$ (\decl{Barbier.reuleauxPiece}, \src{barbier}{Barbier}{Reuleaux}{10}; \decl{Barbier.reuleaux}, \src{barbier}{Barbier}{Reuleaux}{16}).
\end{definition}

On the first branch $R_w$ is the support function of the single corner at distance $w/\sqrt3$ from
the centre; on the second it is the support function of the arc of radius $w$ centred at the
opposite corner. So $h+h''$ is $0$ on the corner pieces and $w$ on the arc pieces, and it jumps at
six points per period.

\begin{theorem}[Reuleaux triangle]\label{thm:reuleaux}
Let $w>0$. Then:
\begin{enumerate}
\item[(a)] $R_w$ is a piecewise constant-width support function of width $w$
(\decl{Barbier.reuleaux_constWidth}, \src{barbier}{Barbier}{Reuleaux}{1074});
\item[(b)] $P(R_w) = \pi w$ (\decl{Barbier.reuleaux_perimeter}, \src{barbier}{Barbier}{Reuleaux}{1083});
\item[(c)] $R_w$ is not twice continuously differentiable (\decl{Barbier.reuleaux_not_smooth}, \src{barbier}{Barbier}{Reuleaux}{1086});
\item[(d)] the points $\gamma(0)$, $\gamma(2\pi/3)$, $\gamma(4\pi/3)$ of the curve of $R_w$ are pairwise
at distance $w$ (\decl{Barbier.reuleaux_vertices}, \src{barbier}{Barbier}{Reuleaux}{1133});
\item[(e)] for every $\theta\in[5\pi/6, 7\pi/6]$, $\lvert\gamma(\theta)-\gamma(0)\rvert = w$
(\decl{Barbier.reuleaux_arc}, \src{barbier}{Barbier}{Reuleaux}{1368});
\item[(f)] the total variation of $\gamma_{R_w}$ on $[0,2\pi]$ is finite
(\decl{Barbier.reuleaux_length_ne_top}, \src{barbier}{Barbier}{ReuleauxLength}{296}), equals our
perimeter (\decl{Barbier.reuleaux_perimeter_eq_eVariationOn}, \src{barbier}{Barbier}{ReuleauxLength}{300}),
and equals $\pi w$ (\decl{Barbier.reuleaux_length}, \src{barbier}{Barbier}{ReuleauxLength}{248}).
\end{enumerate}
\end{theorem}

Part (b) is Theorem~\ref{thm:barbierpw} applied to (a). Part (c) is what makes (b) new: by (c), $R_w$
satisfies no $C^2$ hypothesis, so the smooth Theorem~\ref{thm:barbier} cannot be applied to it. The
proof of (c) computes the one-sided limits of $R_w''$ at $\pi/6$, namely $-w/2$ and $w/2$, and
contradicts continuity of the second derivative. Parts (d) and (e) say the curve really is the
Reuleaux triangle: (d) gives an equilateral triangle of side $w$, and (e) says one side is an arc of
radius $w$ about the opposite corner (the other two follow by the $2\pi/3$-periodicity of $R_w$).

Most of the 1{,}379 lines of \lean{Reuleaux.lean} are in (a): $C^1$-smoothness across the wrap-around
points of the angle reduction, where one-sided derivatives must be matched, and the bookkeeping of
the reduction modulo $2\pi/3$ (Mathlib's \mdecl{toIcoMod}).

\section{The four-vertex theorem}\label{sec:fourvertex}

The \lean{fourvertex} development uses the same curve $\gamma_h$ (\decl{FourVertex.curve}, \src{fourvertex}{FourVertex}{Support}{25}) and the radius of curvature $\rho = h+h''$
(\decl{FourVertex.radius}, \src{fourvertex}{FourVertex}{Support}{17}).

\begin{definition}
A function $g$ has \emph{four sign changes} (\decl{FourVertex.FourSignChanges}, \src{fourvertex}{FourVertex}{Basic}{18}) if there are $t_1<t_2<t_3<t_4<t_1+2\pi$ with $g(t_1)>0$,
$g(t_2)<0$, $g(t_3)>0$, $g(t_4)<0$. A function $f$ has \emph{four vertices}
(\decl{FourVertex.FourVertices}, \src{fourvertex}{FourVertex}{Vertices}{73}) if there are
$c_1<c_2<c_3<c_4<c_1+2\pi$ at which $f$ has, in order, a local maximum, minimum, maximum and minimum.
\end{definition}

\begin{lemma}[Sturm--Hurwitz type]\label{lem:sturm}
Let $g$ be continuous and $2\pi$-periodic with $\int_0^{2\pi}g = \int_0^{2\pi}g\cos =
\int_0^{2\pi}g\sin = 0$. Then $g\equiv 0$ or $g$ has four sign changes
(\decl{FourVertex.eq_zero_or_fourSignChanges}, \src{fourvertex}{FourVertex}{Basic}{150}).
\end{lemma}

\begin{proof}[Proof sketch]
If $g$ has at most two sign changes, choose $a,b$ with $\cos(t-a)-\cos b$ having the same sign
pattern as $g$; then $g(t)(\cos(t-a)-\cos b)\ge 0$ and is not identically $0$, but its integral is
$0$ by the three orthogonality conditions, a contradiction. Sign changes come in pairs by
periodicity.
\end{proof}

\begin{theorem}[four-vertex, radius form]\label{thm:radius}
If $h$ is $C^3$ and $2\pi$-periodic, then $\rho$ is constant or $\rho'$ has four sign changes
(\decl{FourVertex.radius_const_or_fourSignChanges}, \src{fourvertex}{FourVertex}{Support}{124}).
\end{theorem}

Integration by parts shows $\rho'$ is orthogonal to $1$, $\cos$ and $\sin$ (the terms
$h'+h'''$ integrate against $\cos$ and $\sin$ to $0$), so Lemma~\ref{lem:sturm} applies. No convexity
is needed for this analytic statement.

\begin{theorem}[four-vertex theorem for support-function ovals]\label{thm:fourvertex}
Let $h$ be $C^3$ and $2\pi$-periodic with $\rho>0$ everywhere. Then the curvature $1/\rho$ is
constant, or there are $c_1<c_2<c_3<c_4<c_1+2\pi$ at which $1/\rho$ has, in order, a local minimum,
maximum, minimum and maximum (\decl{FourVertex.four_vertex}, \src{fourvertex}{FourVertex}{Vertices}{103}).
\end{theorem}

The step from four sign changes of $\rho'$ to four alternating extrema of $\rho$ is
\decl{FourVertex.fourVertices_of_fourSignChanges}; taking reciprocals swaps maxima and minima. The
oval with support function $5+\cos 2\theta$ is a non-circular example (\decl{FourVertex.oval}, \src{fourvertex}{FourVertex}{Examples}{11}; \decl{FourVertex.oval_four_vertex}, \src{fourvertex}{FourVertex}{Examples}{51}): its radius of curvature ranges over $[2,8]$, attaining
both ends (\decl{FourVertex.bridge_oval}, \src{fourvertex}{FourVertex}{Bridge}{111}).

\section{Bridges to standard notions}\label{sec:bridges}

The theorems above are stated with our own definitions: our perimeter, our constant-width condition,
our curvature. This section proves that those definitions mean what they are meant to.

\subsection{Length}

Mathlib defines the length of a curve on $[a,b]$ as its total variation \mdecl{eVariationOn}: the
supremum over partitions of $\sum\lvert\gamma(t_{k+1})-\gamma(t_k)\rvert$, a value in $[0,\infty]$.
At our Mathlib pin there is no lemma computing it as an integral.

\begin{theorem}[length of a $C^1$ curve]\label{thm:length}
If $\gamma:\RR\to\CC$ is $C^1$ and $a\le b$, then the total variation of $\gamma$ on $[a,b]$ is
finite and equals $\int_a^b\lvert\gamma'(t)\rvert\,dt$
(\decl{Barbier.length_eq_integral_norm_deriv}, \src{barbier}{Barbier}{ArclenFinite}{50}; the
finiteness alone is \decl{Barbier.eVariationOn_ne_top_of_contDiff}, \src{barbier}{Barbier}{ArclenFinite}{7}, and the equation alone
\decl{Barbier.eVariationOn_eq_integral_norm_deriv}, \src{barbier}{Barbier}{Bridge}{144}).
\end{theorem}

\begin{proof}[Proof sketch]
For a partition, $\lvert\gamma(t_{k+1})-\gamma(t_k)\rvert\le\int_{t_k}^{t_{k+1}}\lvert\gamma'\rvert$;
summing gives variation $\le\int_a^b\lvert\gamma'\rvert$, in particular finiteness. For the converse,
let $V(s)$ be the variation on $[a,s]$. Additivity of variation over adjacent intervals and the
two-sided estimate $\lvert\gamma(s+\delta)-\gamma(s)\rvert\le V(s+\delta)-V(s)\le
\int_s^{s+\delta}\lvert\gamma'\rvert$, both $\delta\lvert\gamma'(s)\rvert+o(\delta)$ by uniform
continuity of $\gamma'$, give $V'(s) = \lvert\gamma'(s)\rvert$ from the right, and the fundamental
theorem of calculus concludes.
\end{proof}

The finiteness clause matters in Lean: the conversion of $[0,\infty]$ to $\RR$ sends $\infty$ to $0$,
so the equation alone would not by itself exclude infinite length. It was added after the blind
check of Section~\ref{sec:blind} flagged the difference.

\begin{corollary}[perimeter is length]\label{cor:perlen}
If $h$ is $C^2$ and $2\pi$-periodic with $h+h''\ge 0$, then $P(h)$ is the total variation of
$\gamma_h$ on $[0,2\pi]$ (\decl{Barbier.perimeter_eq_eVariationOn}, \src{barbier}{Barbier}{Bridge}{165}).
\end{corollary}

For the Reuleaux triangle $\gamma'$ is discontinuous, so Corollary~\ref{cor:perlen} does not apply to
it as stated; the Reuleaux triangle is covered separately by Theorem~\ref{thm:reuleaux}(f), and a
general perimeter-equals-length bridge for the whole piecewise class is left open
(Section~\ref{sec:limits}).

\subsection{The convex set}

Mathlib has no support function (its convex-body file lists support sets as future work), so we
construct the set directly.

\begin{definition}
The \emph{body} of $h$ is $K_h = \bigcap_{\theta}\{z : \mathrm{Re}(z\,e^{-i\theta})\le h(\theta)\}$
(\decl{Barbier.body}, \src{barbier}{Barbier}{Bridge}{172}).
\end{definition}

\begin{theorem}[the body]\label{thm:body}
$K_h$ is convex for every $h$ (\decl{Barbier.convex_body}, \src{barbier}{Barbier}{Bridge}{175}). If $h$
is a piecewise constant-width support function of width $w$ then, for every $\theta$:
\begin{enumerate}
\item[(a)] $h(\theta)$ is the greatest value of $\langle\gamma_h(\varphi), e^{i\theta}\rangle$ over
$\varphi\in\RR$, attained at $\varphi = \theta$ (\decl{Barbier.bridge_support_piecewise}, \src{barbier}{Barbier}{BridgePiecewise}{109});
\item[(b)] $\gamma_h(\theta)\in K_h$ (\decl{Barbier.curve_mem_body_piecewise}, \src{barbier}{Barbier}{BridgePiecewise}{170});
\item[(c)] the supremum of $\mathrm{Re}(z\,e^{-i\theta})$ over $z\in K_h$ is $h(\theta)$
(\decl{Barbier.support_body_piecewise}, \src{barbier}{Barbier}{BridgePiecewise}{177});
\item[(d)] the support values of $K_h$ in directions $\theta$ and $\theta+\pi$ add up to $w$
(\decl{Barbier.width_body_piecewise}, \src{barbier}{Barbier}{BridgePiecewise}{184}).
\end{enumerate}
The smooth versions are \decl{Barbier.curve_mem_body}, \decl{Barbier.support_body} and
\decl{Barbier.width_body} (\src{barbier}{Barbier}{Bridge}{180}--197).
\end{theorem}

\begin{theorem}[the body is a convex body]\label{thm:convbody}
If $h$ is a piecewise constant-width support function of width $w$, then $K_h$ is compact
(\decl{Barbier.isCompact_body_piecewise}, \src{barbier}{Barbier}{ConvexBodyBridge}{57}) and nonempty
(\decl{Barbier.body_nonempty_piecewise}, \src{barbier}{Barbier}{ConvexBodyBridge}{60}); with
convexity this makes $K_h$ a Mathlib \mdecl{ConvexBody} whose support values are $h$ and whose width
in every direction is $w$ (\decl{Barbier.exists_convexBody_piecewise}, \src{barbier}{Barbier}{ConvexBodyBridge}{64}). In particular, for $w>0$, the Reuleaux triangle is a Mathlib convex
body of constant width $w$ (\decl{Barbier.reuleaux_convexBody}, \src{barbier}{Barbier}{ConvexBodyBridge}{75}).
\end{theorem}

Boundedness is the only new input: $h$ is continuous and periodic, so $\lvert h\rvert$ is bounded by
$M$; a point $z$ of $K_h$ is bounded by the constraint in the direction of $z$ itself,
$\lVert z\rVert = \langle z, e^{i\arg z}\rangle \le h(\arg z) \le M$.

So $h$ is the support function of the convex set $K_h$, and $K_h$ has width $w$ in every direction:
it is a set of constant width in the standard sense, and by Theorem~\ref{thm:reuleaux} this includes
the Reuleaux triangle. The proof of (a) differentiates
$\varphi\mapsto\langle\gamma_h(\varphi), e^{i\theta}\rangle$, with derivative
$\rho(\varphi)\sin(\theta-\varphi)$ off $S$, and integrates. A supremum of an unbounded set of reals
is $0$ in Lean, so (c) and (d) are proved through attainment (part (a) and part (b)), not read off a
supremum that might be junk.

\subsection{Curvature and simplicity}

Mathlib has no curvature of plane curves, so we define the textbook one.

\begin{definition}
The \emph{signed curvature} of $\gamma:\RR\to\CC$ at $t$ is
$\kappa_\gamma(t) = \mathrm{Im}(\overline{\gamma'(t)}\,\gamma''(t))/\lvert\gamma'(t)\rvert^3$, which is
$(x'y''-y'x'')/(x'^2+y'^2)^{3/2}$ for $\gamma = x+iy$ (\decl{FourVertex.signedCurvature}, \src{fourvertex}{FourVertex}{Curvature}{15}).
\end{definition}

\begin{theorem}[curvature and simplicity]\label{thm:curv}
\begin{enumerate}
\item[(a)] If $h$ is $C^3$ and $\rho(\theta)>0$, then $\kappa_{\gamma_h}(\theta) = 1/\rho(\theta)$
(\decl{FourVertex.signedCurvature_curve}, \src{fourvertex}{FourVertex}{Curvature}{129}).
\item[(b)] If $h$ is $2\pi$-periodic, so is $\gamma_h$ (\decl{FourVertex.curve_periodic}, \src{fourvertex}{FourVertex}{Curvature}{168}).
\item[(c)] If $h$ is $C^3$ and $2\pi$-periodic with $\rho>0$ everywhere, then $\gamma_h$ is injective on $[0,2\pi)$
(\decl{FourVertex.curve_injOn}, \src{fourvertex}{FourVertex}{Curvature}{179}).
\end{enumerate}
\end{theorem}

For (a), $\gamma_h' = \rho\,ie^{i\theta}$ (\decl{FourVertex.bridge_radius}, \src{fourvertex}{FourVertex}{Bridge}{38}) gives $\gamma_h'' = \rho'\,ie^{i\theta}-\rho\,e^{i\theta}$,
so $\overline{\gamma_h'}\gamma_h'' = \rho\rho'+i\rho^2$ and $\kappa = \rho^2/\rho^3$. For (c), the
function $\varphi\mapsto\langle\gamma_h(\varphi),e^{i\theta}\rangle$ has derivative
$\rho(\varphi)\sin(\theta-\varphi)$, so it is strictly increasing on $[\theta-\pi,\theta]$ and strictly
decreasing on $[\theta,\theta+\pi]$, with maximum $h(\theta)$ attained only at $\theta$ in that
window. If $\gamma_h(a) = \gamma_h(b)$ with $a\ne b$ in $[0,2\pi)$, reduce $b$ into the window around
$a$ by periodicity and contradict strictness.

So under the hypotheses of Theorem~\ref{thm:fourvertex}, $\gamma_h$ is a simple closed curve and
$1/\rho$ is its signed curvature. The theorem can therefore be restated with no reference to
support functions in its conclusion:

\begin{theorem}[four-vertex theorem, textbook form]\label{thm:fvtextbook}
Let $h$ be $C^3$ and $2\pi$-periodic with $\rho>0$. Then the signed curvature of the simple closed
curve $\gamma_h$ is constant, or it has, at some $c_1<c_2<c_3<c_4<c_1+2\pi$, a local minimum, maximum,
minimum and maximum, in order (\decl{FourVertex.four_vertex_signedCurvature}, \src{fourvertex}{FourVertex}{Curvature}{228}).
\end{theorem}

The four points are distinct points of the circle $\RR/2\pi\mathbb{Z}$, and each is a local extremum
in Mathlib's sense (\decl{FourVertex.bridge_FourVertices}, \src{fourvertex}{FourVertex}{Bridge}{95};
\decl{FourVertex.bridge_FourSignChanges}, \src{fourvertex}{FourVertex}{Bridge}{64}).

\section{Does the statement mean what it says?}\label{sec:blind}

A complete proof only shows that \emph{the Lean statement} is true. To test whether the Lean
statements say what this paper says, we ran a blind re-formalisation in three rounds.

\paragraph{Round 1 (before this paper).} Fifteen \lean{barbier} results were written out in plain
language by the author's assistant from the code, translated into Lean by an independent model, and
compared. Fourteen pairs were equivalent at once. The fifteenth was the length lemma: the blind
statement asserted finiteness and ours did not. We strengthened ours (Theorem~\ref{thm:length} now
states finiteness), after which that pair was equivalent too.

\paragraph{Round 2 (this paper).} The mathematical sections of this draft, with every code name, source
reference and description of the formalisation removed automatically, were given to a model different
from the one that wrote the proofs. It produced Lean statements of all 22 results and parts of
Sections~\ref{sec:barbier}--\ref{sec:bridges}, reading only that text (self-reported; Mathlib API
lookups allowed). A first attempt was rejected on review because two of its definitions were not
faithful (length defined as the constant~$0$, and the Reuleaux angle reduction omitted); the second
attempt was reviewed definition by definition. A third model then proved, for each pair, that the blind
statement and ours imply each other, using only unfolding of definitions, repackaging and general
Mathlib lemmas, never one of our theorems. The Reuleaux-length and convex-body results were added to
the draft afterwards and checked the same way in a third round, with seven further blind statements
(the only difference being a form one, $\mathrm{ofReal}$ against $\mathrm{toReal}$, in the pair for
Theorem~\ref{thm:reuleaux}(f)).

\paragraph{Result.} Twenty-six of the 29 pairs are equivalent in both directions, including the Reuleaux
support function (which the blind writer defined with a floor function rather than \mdecl{toIcoMod}),
the piecewise hypothesis class, the body and its support values, and both forms of the four-vertex
theorem. A twenty-seventh pair, the Reuleaux perimeter, is equivalent in both directions once the
finiteness of the length (pair 23) is supplied as a separate hypothesis. The other two pairs differed in one direction only, and in the safe one: an earlier version of
this draft stated Theorem~\ref{thm:curv}(a) and~(b) under the hypotheses of part~(c), while the Lean
theorems need fewer (no periodicity for~(a); no smoothness for~(b)). Our statements were therefore
stronger than claimed. We reworded Theorem~\ref{thm:curv} so each part states exactly its own
hypotheses.

\paragraph{Caveats.} The blind writers self-reported which files they read. The comparison establishes
equivalence of statements, not correctness of the prose proofs.

\section{Related work}\label{sec:related}

We searched for prior formal proofs of both theorems in Mathlib (pinned and current), its pull
requests and issues, the Archive of Formal Proofs \cite{afp}, Coq libraries (GeoCoq, MathComp and its analysis
library), HOL Light, HOL4, Mizar, Metamath (\texttt{set.mm} and \texttt{iset.mm}), the ``100
theorems'' and ``1000+ theorems'' lists, GitHub code search, the public Lean Zulip archive, and
several benchmark and autoformalisation collections. We count only a complete, sorry-free proof as
prior art.

\paragraph{Barbier.} We found no formal statement or proof in any system, and no formalisation of the
Reuleaux triangle. The ``1000+'' entry for Barbier's theorem is empty in every system column.

\paragraph{Four-vertex.} The closest work is in \texttt{facebookresearch/atlas-lean} \cite{atlaslean}
(file \texttt{FourVertex.lean}, present from the repository's initial commit of 28 May 2026), which
states an incomplete four-vertex theorem for smooth positive-curvature parametrised simple closed
curves, concluding four stationary points of the curvature rather than the alternating local extrema
used here. Its Sturm--Hurwitz lemma and
its final assembly step are \lean{sorry}. We compiled it at our Mathlib pin (after eight line-local
compatibility edits, no statement changed), and \lean{\#print axioms} confirms that its four-vertex
theorem and its Sturm--Hurwitz lemma depend on \lean{sorryAx}. A second Lean repository
(\texttt{mirajcs/TheFourVertexTheorem}) \cite{mirajcs} has definitions and \lean{sorry}s but no four-vertex
statement. A 2025 Zulip thread records an intention to formalise the theorem, with no code posted. To
our knowledge, Lemma~\ref{lem:sturm} and Theorem~\ref{thm:fourvertex} are the first complete formal
proofs of a Sturm--Hurwitz-type lemma and of a four-vertex theorem.

\paragraph{Support functions.} Mathlib has convex bodies (\mdecl{ConvexBody}) but no support function.
The nearest external Lean work is \texttt{deancureton/MovingSofa} \cite{movingssofa} (September 2026), an
autoformalisation of Baek's \cite{baek2024} moving-sofa proof, which defines support values and directional width of
a set and uses \mdecl{eVariationOn} for boundary length. It states neither theorem treated here.

\paragraph{Our earlier papers.} Two earlier papers in this series formalised classical Euclidean
results by coordinate certificates: Paper~A \cite{paperA} and Paper~B \cite{paperB}. The present paper is the first in the series on
analytic geometry of curves.

\section{Verification}\label{sec:verify}

\paragraph{Gates.} Each repository's gate script is described in Section~\ref{sec:avail}. Both gates passed on two machines of different
architecture, an x86-64 Linux machine and an arm64 macOS machine: \lean{barbier} with 62 declarations
and \lean{fourvertex} with 48, all with only the three standard axioms. The runs were made on
9--10 October 2026 on the sources that were subsequently published; the published commits contain
byte-identical sources (checked by hash), so the gate results apply to them.

\paragraph{Tampering tests.} Each repository has a tampering suite that copies the project, plants one
defect and checks that the gate rejects it at the expected stage: a hidden \lean{sorry}, an
\lean{axiom} declaration split across two lines, a fake \lean{\#print axioms} report printed by
\lean{\#eval}, a macro that could redefine commands, a dropped hypothesis that breaks a proof, and an
extra axiom that only the output parser can catch. All six were rejected in both repositories (run of 10 October 2026, on the x86-64 machine).

\paragraph{Statements not weakened.} The Reuleaux and corner statements were written first, as
statements with \lean{sorry} proofs, then proved by coding agents under the rule that no statement
may change. We compared every statement and definition before and after: all are word-for-word
identical. The pre-existing smooth Barbier and four-vertex theorem statements and proofs are
unchanged by the work of this paper; the bridge modules and the root import files were extended with
new lemmas and imports, while the pre-existing \lean{Barbier/Basic.lean} and \lean{Barbier/Examples.lean}
are byte-for-byte unchanged from their earlier versions.

\paragraph{Statements checked by hand.} The author read every statement quoted in this paper against
the corresponding Lean source by hand, and confirmed that each informal statement agrees with the
formal statement it quotes.

\paragraph{Statement-level checks.} Beyond the gate and the tampering suite, both developments were
run through a statement-level battery. It mutates each theorem (numeric literals, inequalities and
equalities altered) and checks that the proof rejects the mutant; it probes the hypotheses for
consistency; it tries to certify each statement satisfiable; it ablates each hypothesis to see
whether the proof uses it; and it lints for unused variables. Results: \lean{barbier} 386 mutants,
370 rejected, 16 survived; \lean{fourvertex} 91 mutants, 86 rejected, 5 survived. Every survivor is a
mutant that is still \emph{true} --- a looser numeric bound (e.g.\ $2w$ to $3w$), a constant changed in
an auxiliary inequality that stays valid ($\pi/6$ to $\pi/7$), or a strictly stronger hypothesis ---
so survival is expected and is not a defect: the proofs simply do not depend on the exact constant.
The contradiction probe found no contradictory hypothesis set. Satisfiability was certified for 101
of the 168 \lean{barbier} declarations and 26 of the 47 \lean{fourvertex} declarations; the rest are
recorded as unresolved rather than assumed, because no witness was produced. Hypothesis ablation
found every hypothesis used except two in \lean{barbier} (\decl{Barbier.perimeter_eq_eVariationOn}
does not use its convexity hypothesis, and one auxiliary interval lemma does not use its hypothesis
$w>0$), both of which make the theorem slightly more restrictive than necessary without affecting its
correctness.
The unused-variable linter reports no warnings.

\paragraph{Name gate.} A script checks that every declaration named in this paper is in its
repository's gate list and that every source anchor points at a line declaring that name.

\section{Limitations and further work}\label{sec:limits}

\begin{enumerate}
\item \textbf{Barbier for curves given by support functions.} We do not prove that every convex body of
constant width has a support function satisfying Definition~\ref{def:pcw}. That would need the
general theory of support functions of convex bodies (absent from Mathlib) and a regularity argument.
\item \textbf{Four-vertex for ovals.} Theorem~\ref{thm:fvtextbook} covers strictly convex $C^3$ curves
given by $C^3$ support functions with $h+h''>0$. The general theorem for simple closed curves
\cite{kneser1912}, and the converse \cite{dahlberg2005}, are not attempted.
\item \textbf{Perimeter versus length in the piecewise class.} Corollary~\ref{cor:perlen} bridges our
perimeter to Mathlib's length only for smooth support functions, and Theorem~\ref{thm:reuleaux}(f) for
the Reuleaux triangle; we do not prove a general bridge for the whole piecewise class.
\item \textbf{Blind checks} establish equivalence of statements, not correctness of the prose
proofs, and rely on the writers' self-reports (Section~\ref{sec:blind}).
\end{enumerate}

\section*{Acknowledgements and use of AI}
A large-language-model AI assistant was used to write the Lean proofs, to run the prior-art searches
and the verification battery, and to draft this paper. The author directed the work and reviewed the
paper. Every proof is checked by the Lean kernel, and the author takes responsibility for all content.

\begin{thebibliography}{99}

\bibitem{barbier1860}
E.~Barbier,
\emph{Note sur le probl\`eme de l'aiguille et le jeu du joint couvert},
Journal de Math\'ematiques Pures et Appliqu\'ees, 2\textsuperscript{e} s\'erie, tome~5 (1860), 273--286.

\bibitem{mukhopadhyaya1909}
S.~Mukhopadhyaya,
\emph{New methods in the geometry of a plane arc},
Bulletin of the Calcutta Mathematical Society \textbf{1} (1909), 31--37.

\bibitem{kneser1912}
A.~Kneser,
\emph{Bemerkungen \"uber die Anzahl der Extreme der Kr\"ummung auf geschlossenen Kurven und \"uber verwandte Fragen in einer nicht-euklidischen Geometrie},
in: Festschrift Heinrich Weber, Teubner, Leipzig, 1912, 170--180.

\bibitem{gluck1971}
H.~Gluck,
\emph{The converse to the four-vertex theorem},
L'Enseignement Math\'ematique \textbf{17} (1971), 295--309.

\bibitem{osserman1985}
R.~Osserman,
\emph{The four-or-more vertex theorem},
American Mathematical Monthly \textbf{92} (1985), no.~5, 332--337.

\bibitem{dahlberg2005}
B.~E.~J.~Dahlberg,
\emph{The converse of the four vertex theorem},
Proceedings of the American Mathematical Society \textbf{133} (2005), no.~7, 2131--2135.

\bibitem{deturck2007}
D.~DeTurck, H.~Gluck, D.~Pomerleano and D.~S.~Vick,
\emph{The four vertex theorem and its converse},
Notices of the American Mathematical Society \textbf{54} (2007), no.~2, 192--207.

\bibitem{martini2019}
H.~Martini, L.~Montejano and D.~Oliveros,
\emph{Bodies of Constant Width: An Introduction to Convex Geometry with Applications},
Birkh\"auser/Springer, 2019.

\bibitem{lean4}
L.~de Moura and S.~Ullrich,
\emph{The Lean 4 theorem prover and programming language},
in: Automated Deduction -- CADE-28, Lecture Notes in Computer Science \textbf{12699}, Springer, 2021, 625--635.

\bibitem{mathlib2020}
The mathlib Community,
\emph{The Lean mathematical library},
in: Proceedings of the 9th ACM SIGPLAN International Conference on Certified Programs and Proofs (CPP '20), ACM, 2020, 367--381.

\bibitem{afp}
\emph{Archive of Formal Proofs}, ISSN 2150-914x. \url{https://isa-afp.org}

\bibitem{atlaslean}
Meta Fundamental AI Research (FAIR),
\emph{atlas-lean}, GitHub repository, 2026. \url{https://github.com/facebookresearch/atlas-lean}

\bibitem{mirajcs}
\emph{mirajcs/TheFourVertexTheorem},
GitHub repository, 2026. \url{https://github.com/mirajcs/TheFourVertexTheorem}

\bibitem{movingssofa}
D.~Cureton,
\emph{MovingSofa} (Lean formalisation of Baek's moving-sofa proof),
GitHub repository, 2026. \url{https://github.com/deancureton/MovingSofa}

\bibitem{baek2024}
J.~Baek,
\emph{Optimality of Gerver's sofa}, arXiv:2411.19826 [math.MG], 2024.

\bibitem{paperA}
J.~Bald,
\emph{Classical plane geometry by coordinate certificates: Desargues, Monge, Simson, Napoleon and Lami in Lean~4},
Zenodo, 2026. \url{https://doi.org/10.5281/zenodo.23239930}

\bibitem{paperB}
J.~Bald,
\emph{The Malfatti circles in Lean~4, with independent certificate proofs of Feuerbach, Miquel and Brahmagupta},
Zenodo, 2026. \url{https://doi.org/10.5281/zenodo.23240500}

\end{thebibliography}

\end{document}
